\documentclass[11pt]{article}
\usepackage{fontspec}
\setmainfont{Times New Roman}
\setsansfont{Arial}
\setmonofont{Consolas}
\usepackage{amsmath,amssymb,mathtools}
\usepackage{geometry}
\usepackage{microtype}
\geometry{a4paper,margin=2.35cm}
\setlength{\parindent}{1.25em}
\setlength{\parskip}{0pt}
\makeatletter
\renewcommand\footnotesize{\@setfontsize\footnotesize{7.7}{8.7}\selectfont}
\makeatother
\emergencystretch=75em
\allowdisplaybreaks
\sloppy
\hbadness=10000
\vbadness=10000
\hfuzz=2pt
\vfuzz=10pt
\raggedbottom

\begin{document}

% Унит: Папер31 Сецтион02 ентры 1 в001. Дирецт Интерславиц Латин транслатион фром the Герман финал-аудитед соурце слице, соурце линес 92--100 / сегментс P31-С0028--P31-С0032. Фоотноте цоунтер цонтинуес афтер Папер31 Сецтион01 ентры 3.
\setcounter{footnote}{16}

\subsection*{\S 2. Колца конечного ранга. Представјенье како директна сума}

\textbf{1. Колца конечного ранга.} Ако $\mathfrak R$ јест конечны $P$-модул, тогда $\mathfrak R$ имаје конечны ранг, рецимо $k$, взходно к $P$, т. ј. в $\mathfrak R$ сут $k$ и не веч линеарно независных взходно к $P$ елементов, и всакы систем од $k$ линеарно независных взходно к $P$ елементов из $\mathfrak R$ зато твори модулну базу $\mathfrak R$. Ако $P$-модул $\mathfrak B$ конечного ранга имаје како делитель другы $\mathfrak A$ того самого конечного ранга, $\mathfrak B=O(\mathfrak A)$, тогда оны сут идентичне.

Одтепер $\mathfrak R$ предполагаје се како конечного ранга $n$ взходно к $P$. Тогда всакы $P$-модул из $\mathfrak R$ имаје конечны ранг; зато в всакој властној цепи делительев или кратников $P$-модулов из $\mathfrak R$ ранг всакого модула мора быти вечши, односно менши, неж ранг непосреднье предходного; цепи прекидајут се по конечно много кроках. Особливо теды важи "теорем о двојној цепи" за систем всих идеалов из $\mathfrak R$.

Из того следује,\footnote{Idealtheorie, \S 7, 2.} же просты идеал $\mathfrak p$ из $\mathfrak R$ не имаје никакого властного делительа разного од јединичного идеала; зато колцо класов остатков $\mathfrak R/\mathfrak p$ по простом идеалу разном од јединичного става полье. Ако $\mathfrak q$ јест примарны идеал принадлежечи к $\mathfrak p$, т. ј. $\mathfrak q=O(\mathfrak p)$ и $\mathfrak p^{\rho}=O(\mathfrak q)$, тогда колцо класов остатков $\mathfrak R/\mathfrak q$ става примарно колцо, т. ј. колцо, в ктором некоја степень (ту во всаком случају уже $\rho$-та) всакого нулового делительа изчезаје, и $\mathfrak R/\mathfrak q$ садржи само нулове делителье и јединице. Последнье следује непосреднье из того, же всакы идеал не делимы чрез $\mathfrak p$ става взајемно просты с $\mathfrak p$, и зато такоже с $\mathfrak q$.

Из теорема о двојној цепи и из ексистенције јединичного елемента далеј следује,\footnote{Idealtheorie, \S 7, 3. Теорем V.} же всакы идеал $\mathfrak a$ из $\mathfrak R$ може једнозначно быти представјен како продукт конечно многих по парах взајемно простых примарных идеалов. Чрез взајемну простоту тој продукт става равны најменшему обчему кратнику и одговарја представјеньу колца класов остатков $\mathfrak R/\mathfrak a$ како директној сумы примарных колц,\footnote{Idealtheorie, \S 4, 5. То, что јест додано в нынешнем своду, находи се много где разсејано в литературе.} т. ј. $\mathfrak R/\mathfrak a$ става највечи обчи делитель тых колц, и представјенье всакого елемента из $\mathfrak R/\mathfrak a$ како сумы јест једнозначно.

\end{document}

